\documentclass[a4paper,12pt]{article}

%% ============================================
%% 经济学论文 LaTeX 模板
%% 基于 AER (American Economic Review) 格式
%% ============================================

% 页面设置
\usepackage[top=1in,bottom=1in,left=1in,right=1in]{geometry}
\usepackage{setspace}
\doublespacing

% 编码
\usepackage[utf8]{inputenc}
\usepackage[T1]{fontenc}
\usepackage{indentfirst}
\usepackage{microtype}

% 数学
\usepackage{amsmath, amssymb, amsthm}
\numberwithin{equation}{section}

% 表格与图表
\usepackage{booktabs}
\usepackage{graphicx}
\usepackage{longtable}
\usepackage{array}
\usepackage{multirow}

% 交叉引用
\usepackage{hyperref}
\hypersetup{
    colorlinks=true,
    linkcolor=blue,
    citecolor=blue,
    urlcolor=blue
}

% 参考文献格式
\usepackage{natbib}
\bibliographystyle{aea}

% 自定义命令
\usepackage{xargs}
\usepackage{ifthen}

% 定义实证经济学常用命令
\newcommand{\E}{\mathbb{E}}
\newcommand{\Var}{\text{Var}}
\newcommand{\Cov}{\text{Cov}}
\newcommand{\indep}{\perp\!\!\!\perp}
\DeclareMathOperator{\plim}{plim}
\DeclareMathOperator{\argmax}{argmax}
\DeclareMathOperator{\argmin}{argmin}

% 表格小数对齐
\newcolumntype{d}[1]{D{.}{.}{#1}}
\newcolumntype{.}{>{\current@color}}
\usepackage{etoolbox}
\AtBeginEnvironment{tabular}{\sffamily}

% 论文标题页
\title{\textbf{[标题]}\\}
\author{
\textbf{[作者姓名]}\\[1ex]
\textit{[机构]}\\[1ex]
\textit{[邮箱]}\\[2ex]
\textbf{\today}
}
\date{}

\begin{document}
\maketitle
\thispagestyle{empty}

% ============================================
% 摘要
% ============================================
\vspace{-2ex}
\begin{center}
\textbf{ABSTRACT}
\end{center}
\vspace{-1ex}

\noindent
[在此处撰写摘要，150-300字。摘要应简洁明了，包含以下要素：]

\begin{itemize}
    \item \textbf{研究问题}: 本文研究...
    \item \textbf{方法}: 本文使用...方法，利用...数据...
    \item \textbf{主要发现}: 研究发现...
    \item \textbf{结论}: 政策含义是...
\end{itemize}

\vspace{1ex}
\noindent
\textit{Keywords:} [关键词1, 关键词2, 关键词3] \\
\textit{JEL Classification:} [JEL代码，如：O10, O40, C50]

\newpage

% ============================================
% 1. 引言
% ============================================
\section{Introduction}
\label{sec:intro}

\subsection{研究背景}
[描述研究主题的现实背景和重要性]

\subsection{研究问题}
[明确提出本文要回答的核心问题]

\subsection{研究意义}
[理论贡献和实践价值]

\subsection{研究创新}
[本文与现有文献相比的主要创新点]

\subsection{文章结构}
[简要说明后续各部分内容]

% ============================================
% 2. 文献综述
% ============================================
\section{Literature Review}
\label{sec:lit}

\subsection{核心概念界定}
[定义本文使用的核心概念]

\subsection{理论框架}
[理论基础和研究假设]

\subsection{相关文献回顾}
[从以下角度回顾文献：]

\subsubsection{[主题1]}
[相关研究概述，包括作者(年份)的研究发现]

\subsubsection{[主题2]}
[相关研究概述]

\subsection{文献评述与研究缺口}
[总结现有研究的不足，说明本文如何填补]

% ============================================
% 3. 数据与研究方法
% ============================================
\section{Data and Methodology}
\label{sec:data}

\subsection{数据来源}
[描述数据来源、时间跨度、样本范围]

\subsection{变量定义}
[使用表格清晰定义所有变量]

\begin{table}[htbp]
\caption{变量定义表}
\label{tab:var_def}
\centering
\begin{tabular}{lcl}
\toprule
变量名称 & 符号 & 定义 \\
\midrule
被解释变量 & $Y_{it}$ & [定义] \\
解释变量 & $X_{it}$ & [定义] \\
控制变量 & $Z_{it}$ & [定义] \\
\bottomrule
\end{tabular}
\end{table}

\subsection{描述性统计}
[报告样本基本特征]

\subsection{模型设定}
[设定基准回归模型]

\begin{equation}
\label{eq:baseline}
Y_{it} = \alpha + \beta X_{it} + \gamma Z_{it} + \mu_i + \lambda_t + \varepsilon_{it}
\end{equation}

其中，$\mu_i$ 和 $\lambda_t$ 分别表示国家和年份固定效应。

\subsection{识别策略}
[说明因果识别的假设和策略]

\subsubsection{内生性处理}
[讨论潜在的内生性来源及处理方法]

\subsubsection{稳健性检验}
[说明将进行的稳健性检验]

% ============================================
% 4. 实证结果
% ============================================
\section{Empirical Results}
\label{sec:results}

\subsection{基准回归}
[报告基准回归结果]

\begin{table}[htbp]
\caption{基准回归结果}
\label{tab:baseline}
\centering
\begin{tabular}{lccc}
\toprule
 & \multicolumn{3}{c}{被解释变量 $Y$} \\
\cmidrule(lr){2-4}
 & (1) & (2) & (3) \\
\midrule
解释变量 $X$ & 0.123^{***} & 0.098^{**} & 0.087^{*} \\
            & (0.032) & (0.041) & (0.045) \\
控制变量 $Z$ & & 0.215^{*} & 0.198 \\
            & & (0.112) & (0.123) \\
\midrule
固定效应 & 否 & 国家 & 双向 \\
观测值 & 1,234 & 1,234 & 1,234 \\
$R^2$ & 0.123 & 0.245 & 0.389 \\
\bottomrule
\end{tabular}
\begin{minipage}{\textwidth}
\smallskip
\textit{注:} 括号内为稳健标准误，$^{*}$、$^{**}$、$^{***}$分别表示10\%、5\%、1\%显著性水平。
\end{minipage}
\end{table}

\subsection{稳健性检验}
[报告各种稳健性检验结果]

\subsubsection{替换核心变量}
[使用替代变量进行检验]

\subsubsection{子样本分析}
[按不同子样本进行分析]

\subsubsection{工具变量法}
[处理内生性的结果]

\subsection{异质性分析}
[按不同群体分析结果的异质性]

% ============================================
% 5. 结论
% ============================================
\section{Conclusion}
\label{sec:concl}

\subsection{主要发现}
[总结本文主要实证发现]

\subsection{政策建议}
[基于研究发现提出政策建议]

\subsection{研究局限}
[承认研究的局限性]

\subsection{未来研究方向}
[提出未来可拓展的方向]

% ============================================
% 参考文献
% =========================================---
\newpage
\bibliography{references}

% ============================================
% 附录
% =========================================---
\newpage
\appendix
\section*{附录}

\subsection{表A1：变量描述性统计}
% [完整的描述性统计表]

\subsection{表A2：稳健性检验结果}
% [额外的稳健性检验]

\subsection{表A3：工具变量检验}
% [IV相关检验结果]

\end{document}